We design three domain-specific guidance heuristics ($h_1$, $h_2$, $h_4$) and one search-control strategy ($h_3$) to guide A* toward reliable, low-risk, and realistically executable arbitrage routes. Each heuristic captures a different dimension of execution feasibility. Rather than detecting theoretical profit first and evaluating risk afterward, our approach embeds execution awareness directly into the search process.

Each node-level heuristic modifies the evaluation function
\[
f(n) = g(n) + h(n),
\]
where $g(n)$ is the accumulated cost in log space and $h(n) \geq 0$ is an additive guidance penalty encoding domain knowledge about future execution feasibility. As discussed in Section~3, these heuristics are \emph{guidance penalties}: they do not modify $g(n)$ and do not claim admissibility.

\subsection{$h_1$ -- Liquidity Heuristic}

The liquidity heuristic models the risk that low-volume markets cannot support the required trade size without excessive execution impact.

Let
\begin{align}
\textit{vol\_per\_sec}
&=
\frac{\textit{24h\_quote\_volume}}{24 \times 3600}
\\
\textit{liquidity\_ratio}
&=
\frac{\textit{vol\_per\_sec} \cdot T_{\textit{remain}}}
{\textit{order\_size\_usd}}.
\end{align}

This ratio estimates whether expected market flow during the remaining arbitrage window can support the intended order size.

We normalize this quantity into a bounded score:

\[
\textit{score}_{liq} =
\textit{clamp}\!\left(
\frac{\log_{10}(\textit{liquidity\_ratio}) + 1}{2},
0,
1
\right).
\]

The heuristic cost becomes:

\[
h_1(n) = \lambda_{liq} \big(1 - \textit{score}_{liq}(n)\big).
\]

\textbf{Characteristics:}
\begin{itemize}
    \item Time-aware liquidity modeling
    \item Accounts for order size relative to trading activity
    \item Favors deep, actively traded markets
\end{itemize}

This heuristic is well suited for conservative strategies prioritizing reliable execution over aggressive profit maximization.

\subsection{$h_2$ -- Slippage Heuristic}

The slippage heuristic directly models expected execution price impact using real-time order-book depth.

Let $P_{mid}$ denote the mid-price and $\textit{VWAP}(Q)$ the volume-weighted average price required to fill order size $Q$. Slippage in basis points is defined as:

\[
\textit{slippage}_{bps} =
\frac{\textit{VWAP}(Q) - P_{mid}}{P_{mid}} \times 10^{4}.
\]

The heuristic is defined as:

\[
h_2(n) =
\lambda_{slip} \cdot 
\max \left\{ 0, \textit{slippage}_{bps} - \theta \right\},
\]

where $\theta$ is a configurable slippage tolerance (default: 10 bps = 0.10\%).

\paragraph{Relationship to edge costs (no double-counting).}
Edge effective rates $r(e)$ encode \emph{static} costs: the posted taker fee (e.g., 0.10\% on Binance) and withdrawal/gas fees. These fees are deterministic and independent of order size. By contrast, $h_2$ estimates \emph{dynamic} order-book slippage: the price impact of walking the live order book with a specific order size $Q$. Because edges use fixed taker-fee rates rather than order-book-derived VWAP, $h_2$ is additive guidance that captures information not present in $g(n)$, and there is no double-counting.

\textbf{Characteristics:}
\begin{itemize}
    \item Uses real-time order book data
    \item Side-aware (buy vs. sell impact)
    \item Directly penalizes large price impact
    \item Purely additive guidance: does \emph{not} modify $g(n)$
\end{itemize}

This heuristic performs particularly well when order-book depth is the dominant constraint.

\subsection{Multi-Start Search Strategy ($h_3$)}
\label{sec:h3_multistart}

Unlike $h_1$, $h_2$, and $h_4$, which are node-level heuristic functions that return a scalar penalty for each node, $h_3$ is a \textbf{search-control strategy} that operates at the meta-level. It runs multiple independent A* searches concurrently from randomly selected starting nodes, each using a base heuristic (typically $h_1$), and returns the best solution found across all parallel instances.

This strategy is \emph{not} a heuristic in the formal A* sense: it does not define a per-node cost estimate $h_3(n)$, and it does not affect the $f(n) = g(n) + h(n)$ evaluation within any single search. Instead, it addresses a fundamentally different problem: when the profitable region of the graph is unknown a priori, multi-start search diversifies exploration across the state space.

Formally, let $\{v_1, v_2, \dots, v_k\}$ be $k$ randomly sampled starting nodes. For each $v_i$, we run an independent A* search with some base heuristic $h_b \in \{h_1, h_2, h_4\}$:
\[
\text{result}_i = A^*(v_i, h_b), \quad i = 1, \dots, k.
\]
The final output is $\arg\max_{i} \text{profit}(\text{result}_i)$ among successful searches.

\textbf{Characteristics:}
\begin{itemize}
    \item Diversifies search across starting exchanges
    \item Mitigates dependence on a fixed starting point
    \item Embarrassingly parallel (scales with available threads)
    \item Higher total computational cost (linear in $k$)
\end{itemize}

This approach is particularly useful when profitable starting points are not known in advance.

\subsection{$h_4$ -- Chain Congestion and Exchange Reliability Heuristic}

The $h_4$ heuristic incorporates transfer latency and operational reliability into the search.

Let $t_{min}$ denote the shortest outgoing transfer time from the current node, and $T_{remain}$ the remaining arbitrage window. The chain reliability score is:

\[
\textit{score}_{chain} =
1 - \frac{t_{min}}{T_{remain}}.
\]

The corresponding chain penalty is:

\[
h_{chain}(n) =
\lambda_{chain} \big(1 - \textit{score}_{chain}(n)\big).
\]

If each exchange is assigned a reliability score $s(n) \in [0,1]$, we define:

\[
h_{exchange}(n) =
\lambda_{exchange} \big(1 - s(n)\big).
\]

The combined heuristic is:

\[
h_4(n) = h_{chain}(n) + h_{exchange}(n).
\]

\paragraph{Chain transfer times.}
Transfer-time estimates $t_{\min}$ are drawn from documented blockchain finality times plus typical exchange processing delays. For example: Solana $\approx$~1\,s, BNB Chain $\approx$~4\,s, Polygon $\approx$~5\,s, Tron $\approx$~30\,s, Arbitrum/Base $\approx$~120\,s, Ethereum $\approx$~600\,s.

\paragraph{Weighted A* with $h_4$.}
As described in Section~3, $h_4$ is used with Weighted A*. The node-specific weight $w(n)$ amplifies the heuristic penalty for nodes on fast-transfer chains where arbitrage opportunities are fleeting.

\textbf{Characteristics:}
\begin{itemize}
    \item Models blockchain congestion and finality risk
    \item Incorporates exchange operational reliability
    \item Uses Weighted A* ($w \in [1.0, 3.0]$) for risk-adaptive search
    \item Prioritizes fast and stable execution paths
\end{itemize}